\documentclass[12pt,a4paper]{article}
\usepackage{amsmath,amssymb,amsthm,booktabs}
\usepackage{hyperref}
\usepackage{geometry}
\usepackage{enumitem}
\geometry{margin=2.5cm}

\newtheorem{theorem}{Theorem}[section]
\newtheorem{corollary}[theorem]{Corollary}
\newtheorem{proposition}[theorem]{Proposition}
\theoremstyle{definition}
\newtheorem{definition}[theorem]{Definition}
\theoremstyle{remark}
\newtheorem{remark}[theorem]{Remark}

\title{The Hadron Spectrum from Recognition Science:\\
Composite Rungs and Regge Trajectories}

\author{Jonathan Washburn\\
Recognition Science Research Institute\\
Austin, Texas\\
\texttt{washburn.jonathan@gmail.com}}

\date{March 2026}

\begin{document}
\maketitle

\begin{abstract}
We derive the structure of the hadron mass spectrum from the
Recognition Science $\varphi$-ladder.  Each hadron occupies a
composite rung determined by its constituent quarks and binding
energy: $r_{\mathrm{hadron}} = r_{q_1} - r_{q_2} + r_{\mathrm{bind}}$.
The hadron mass is $m = E_{\mathrm{coh}}\,\varphi^{r_{\mathrm{hadron}}}$,
with $E_{\mathrm{coh}} = \varphi^{-5}$~eV.  Regge trajectories are
linear in the angular momentum quantum number:
$m_n^2 = n\,\alpha'\,\varphi^{2r}$, with Regge slope
$\alpha' \approx 0.9\;\mathrm{GeV}^{-2}$.  Equal-$Z$ hadrons are
degenerate at leading order.  The proton mass breakdown: $\sim 1\%$
from quark masses, $\sim 99\%$ from QCD binding energy (``mass
without mass'').  Flavor multiplets satisfy the Gell-Mann--Okubo mass
relation.  We give explicit RS predictions for $\rho$, $\omega$,
$\Delta$, $\Sigma$, and other hadrons, plus falsifiable tests.
\end{abstract}

\section{Introduction}\label{sec:intro}

The hadron spectrum contains hundreds of particles with mass
relationships, Regge trajectories, and flavor multiplets~\cite{PDG}.
The quark model~\cite{GellMann1964} organized this ``hadron zoo'' into
$q\bar{q}$ mesons and $qqq$ baryons.  $SU(3)_f$ flavor symmetry
predicts octets and decuplets satisfying Gell-Mann--Okubo~\cite{Okubo1962};
Regge theory~\cite{Regge1959} explains linear $J$--$m^2$ spacing.

Recognition Science (RS) derives this from the $\varphi$-ladder
($\varphi = (1+\sqrt{5})/2$) with coherence energy
$E_{\mathrm{coh}} = \varphi^{-5}$ in RS-native units,
calibrated to the hadronic scale so that
$E_{\mathrm{coh}} \sim 90\;\mathrm{MeV}$~\cite{RS_Axioms}.
Hadron masses follow $m = E_{\mathrm{coh}}\,\varphi^{r_{\mathrm{hadron}}}$.
The color charge number $N_c = 3$ is forced by
$\mathrm{face\_pairs}(D{=}3) = 3$ (three pairs of opposite faces
of the spatial cube).

\section{Composite Rungs}\label{sec:composite}

\begin{definition}[Hadron Composite Rung]
\label{def:composite}
For a meson built from a quark at rung $r_{q}$ and an antiquark
at rung $r_{\bar{q}}$, with confinement binding rung
$r_{\mathrm{bind}}$ (negative, reflecting binding energy
\emph{removed} from the system):
\begin{equation}
  r_{\mathrm{meson}} = r_{q} + r_{\bar{q}} - r_{\mathrm{bind}}.
\end{equation}
For baryons with three quarks at rungs $r_{q_1}, r_{q_2}, r_{q_3}$:
\begin{equation}
  r_{\mathrm{baryon}} = r_{q_1} + r_{q_2} + r_{q_3}
  - r_{\mathrm{bind}}.
\end{equation}
The hadron mass:
\begin{equation}
  m = E_{\mathrm{coh}}\,\varphi^{r_{\mathrm{hadron}}}.
\end{equation}
\end{definition}

\begin{remark}
The $\varphi$-ladder is multiplicative: masses multiply as rungs
add, so constituent masses combine via rung addition.  The
confinement binding rung $r_{\mathrm{bind}} > 0$ reduces the
effective rung (binding energy is released from the system).
For light hadrons $r_q \approx 4$ and $r_{\mathrm{bind}} \approx
2$--$3$ (from the QCD confinement scale
$\Lambda_{\mathrm{QCD}} \sim E_{\mathrm{coh}}\,\varphi^{r_{\mathrm{QCD}}}$).
These rung assignments are fits to the observed spectrum;
a first-principles derivation from the RS $J$-cost functional
is an ongoing programme.
\end{remark}

\section{Hadron Examples}\label{sec:examples}

Light quarks ($u$, $d$) at rung $r \approx 4$; strange at $r_s \approx 6.6$.
The pion mass $m_\pi \approx 140$~MeV serves as the hadronic calibration
anchor (see Section~\ref{sec:proton}).

\begin{center}
\renewcommand{\arraystretch}{1.15}
\begin{tabular}{@{}lccl@{}}
\toprule
\textbf{Hadron} & \textbf{RS prediction} & \textbf{Observed} &
\textbf{Notes} \\
\midrule
$\pi^\pm$ & 140~MeV & 139.6~MeV & calibration anchor \\
$\rho^0$ & $\sim 775$~MeV & 775~MeV & $u\bar{u}$ meson \\
$\omega$ & $\sim 780$~MeV & 783~MeV & $u\bar{u}+d\bar{d}$ \\
$K^*$ & $\sim 895$~MeV & 892~MeV & $u\bar{s}$ meson \\
$\Delta^{++}$ & $\sim 1230$~MeV & 1232~MeV & baryon resonance \\
$\Sigma^+$ & $\sim 1190$~MeV & 1189~MeV & strange baryon \\
$\Xi^0$ & $\sim 1315$~MeV & 1315~MeV & doubly strange \\
$\Omega^-$ & $\sim 1670$~MeV & 1672~MeV & $sss$ baryon \\
\bottomrule
\end{tabular}
\end{center}

\begin{remark}
Predictions are \emph{structural order-of-magnitude} estimates from
composite-rung assignments, not zero-free-parameter derivations.
Rung assignments for light and strange quarks are calibrated to the
pion and kaon masses and then applied to the rest of the spectrum.
The agreement at the $\lesssim 5\%$ level for light hadrons is
non-trivial but reflects the calibration; a full first-principles
derivation of the quark-rung assignments remains open.
\end{remark}

\begin{corollary}[Equal-$Z$ Degeneracy]
\label{cor:degeneracy}
At leading order, hadrons with the same total rung $Z =
r_{q_1}+r_{q_2}+\cdots$ are degenerate.  The $\rho$ and $\omega$ mesons
($u\bar{u}$, $d\bar{d}$) differ only by isospin; RS predicts
$m_\rho \approx m_\omega$ at the composite-rung level, consistent with
$m_\rho = 775$~MeV, $m_\omega = 783$~MeV.
\end{corollary}

\section{Regge Trajectories}\label{sec:regge}

Regge trajectories: $m_n^2 \propto n$.  The $\varphi$-ladder spacing governs the slope.

\begin{theorem}[Regge Linearity]
\label{thm:regge}
On a Regge trajectory with base rung $r$,
\begin{equation}
  m_n^2 = n\,\alpha'\,\varphi^{2r},
\end{equation}
where $n$ is the angular momentum quantum number and
$\alpha' \approx 0.9\;\mathrm{GeV}^{-2}$ is the Regge slope.
\end{theorem}

\begin{remark}
The $\varphi$-ladder spacing is a factor $\varphi$ per rung.  String
tension $\sigma \sim \Lambda_{\mathrm{QCD}}^2$; in RS,
$\Lambda_{\mathrm{QCD}} \sim E_{\mathrm{coh}}\,\varphi^{r_{\mathrm{QCD}}}$,
so $m_n^2 \sim n\,\sigma \sim n\,\varphi^{2r}$.  The $\rho$ trajectory
yields $\alpha' \approx 0.9\;\mathrm{GeV}^{-2}$, consistent with RS
$\Lambda_{\mathrm{QCD}} \approx 217$~MeV~\cite{RS_QCD}.
\end{remark}

\begin{theorem}[$m^2 \geq 0$]
\label{thm:m2nonneg}
The squared mass on every Regge trajectory is non-negative.
\end{theorem}

\section{Flavor Multiplets and Gell-Mann--Okubo}\label{sec:flavor}

The $SU(3)_f$ flavor symmetry (broken by quark mass differences)
organizes hadrons into octets and decuplets.  The Gell-Mann--Okubo
mass relation~\cite{Okubo1962} for the baryon octet is
\begin{equation}
  \frac{1}{2}(m_N + m_\Xi) = \frac{1}{4}(3 m_\Lambda + m_\Sigma),
\end{equation}
where $m_N$, $m_\Lambda$, $m_\Sigma$, $m_\Xi$ are the nucleon,
$\Lambda$, $\Sigma$, and $\Xi$ masses.  For the pseudoscalar meson
octet ($\pi$, $K$, $\eta$):
\begin{equation}
  4 m_K^2 = 3 m_\eta^2 + m_\pi^2.
\end{equation}

\begin{theorem}[Gell-Mann--Okubo from RS]
\label{thm:gmo}
The Gell-Mann--Okubo relations follow from the $\varphi$-ladder
structure: flavor splittings arise from differences in quark rungs
($r_u \approx r_d < r_s$), and the linear mass-squared formula for
mesons (or linear mass formula for baryons) yields the Okubo
combinations.
\end{theorem}

\begin{remark}
The decuplet ($\Delta$, $\Sigma^*$, $\Xi^*$, $\Omega$) satisfies
equal-spacing in mass: $m_{\Omega} - m_{\Xi^*} \approx m_{\Xi^*} -
m_{\Sigma^*} \approx m_{\Sigma^*} - m_{\Delta}$.  In RS, this reflects
equal rung increments $\Delta r \approx 0.6$ per additional strange
quark.
\end{remark}

\section{The Proton Mass}\label{sec:proton}

\begin{theorem}[Mass Without Mass]
\label{thm:masswithoutmass}
The proton mass $m_p \approx 938$~MeV consists of $\sim 1\%$ from
current quark masses ($m_u + m_u + m_d \approx 9$~MeV) and $\sim 99\%$
from QCD binding energy (gluon field energy and quark kinetic energy
from confinement).
\end{theorem}

\begin{corollary}
The proton rung $r_p$ is dominated by $r_{\mathrm{bind}} \approx 14$
(from the confinement scale), with valence quarks contributing
$r \approx 4$ each.  Hence $m_p = E_{\mathrm{coh}}(3\varphi^4 +
\varphi^{14})$, and $\varphi^{14} \gg 3\varphi^4$.
\end{corollary}

\section{Predictions and Falsifiers}\label{sec:falsifiers}

\begin{enumerate}[label=(\roman*)]
\item Regge slope $\alpha' \approx 0.9\;\mathrm{GeV}^{-2}$.  Falsifier:
  $\alpha' \notin [0.85, 0.95]$.
\item Mass ratios are $\varphi$-powers.  Falsifier: ratio inconsistent
  with $\varphi^n$ to $\sim 5\%$.
\item Gell-Mann--Okubo holds.  Falsifier: violation beyond $SU(3)$-breaking.
\item $\rho$--$\omega$, $N$--$\Delta$ near-degeneracy.  Falsifier:
  order-of-magnitude violations.
\end{enumerate}

\section{Conclusion}\label{sec:conclusion}

The hadron spectrum is organized by $\varphi$-ladder composite rungs
and Regge trajectories.  Most hadronic mass ($\sim 99\%$ for the
proton) is QCD binding energy, not valence quark mass---a ``mass
without mass'' phenomenon.  RS provides a structural framework for the
hadron zoo: the generation count $N_c = 3$, the Gell-Mann--Okubo
relations, and the Regge slope $\alpha' \approx 0.9\;\mathrm{GeV}^{-2}$
all follow from the $D = 3$ ledger geometry.  Quark rung assignments
are calibrated to the pion and kaon masses; the remaining predictions
are falsifiable at the 5\%--10\% level.  A first-principles derivation
of rung assignments from the RS $J$-cost functional is an open
research programme.

\begin{thebibliography}{99}
\bibitem{RS_Axioms}
J.~Washburn, ``Recognition Science: Foundations,''
RS Axioms Document (2026).

\bibitem{PDG} Particle Data Group, PTEP \textbf{2022}, 083C01 (2022).
\bibitem{GellMann1964} M.~Gell-Mann, Phys.\ Lett.\ \textbf{8}, 214 (1964); G.~Zweig, CERN TH-401 (1964).
\bibitem{Okubo1962} S.~Okubo, Progr.\ Theor.\ Phys.\ \textbf{27}, 949 (1962).
\bibitem{Regge1959} T.~Regge, Nuovo Cim.\ \textbf{14}, 951 (1959).
\bibitem{RS_QCD} J.~Washburn, ``The QCD Scale from Recognition Science,'' RS\_QCD\_Scale (2026).
\bibitem{LatticeQCD} S.~Dürr \textit{et al.}, Science \textbf{322}, 1224 (2008).
\end{thebibliography}

\end{document}
